\chapter{Boolean Analysis — Torus}
%%% generated by the blueprint pipeline from TCSlib.BooleanAnalysis.GaussianSpace.Noise.Torus
\section{Overview}

Declarations of the Lean module \texttt{TCSlib.BooleanAnalysis.GaussianSpace.Noise.Torus}.

\section{Declarations}

\begin{definition}[Gaussian torus map]
\lean{BooleanAnalysis.GaussianSpace.gaussianTorusMap}
\label{BooleanAnalysis.GaussianSpace.gaussianTorusMap}
\leanok
For a natural number $n$ and a vector $x=(x_i)_{i=0}^{n-1}\in\bbr^n$, define a point of
the torus with coordinates indexed by pairs $(i,\beta)\in\mathbin
n\times\\{\mathrm{false},\mathrm{true}\}$: its $(i,\mathrm{false})$ coordinate is $x_i$
modulo integers, and its $(i,\mathrm{true})$ coordinate is $\sqrt{2}\,x_i$ modulo
integers.
\end{definition}

\begin{lemma}[Injectivity of the Gaussian torus map]
\lean{BooleanAnalysis.GaussianSpace.gaussianTorusMap_injective}
\label{BooleanAnalysis.GaussianSpace.gaussianTorusMap_injective}
\leanok
\uses{BooleanAnalysis.GaussianSpace.gaussianTorusMap}
\difficulty{6}
For every nonnegative integer $n$, the map sending $x\in\mathbb R^n$ to the
product-torus point whose coordinates are $x_i$ and $\sqrt{2}x_i$ modulo $1$ is
injective.
\end{lemma}

\begin{lemma}[Continuity of the Gaussian torus map]
\lean{BooleanAnalysis.GaussianSpace.continuous_gaussianTorusMap}
\label{BooleanAnalysis.GaussianSpace.continuous_gaussianTorusMap}
\leanok
\uses{BooleanAnalysis.GaussianSpace.gaussianTorusMap}
\difficulty{4}
For each natural number $n$, the map from $\mathbb R^n$ to the product torus indexed by
$\{1,\ldots,n\}\times\{0,1\}$ that sends $x$ to the coordinates represented by $x_i$ and
$\sqrt{2}x_i$ modulo $1$ is continuous.
\end{lemma}

\begin{lemma}[Measurable embedding into the torus]
\lean{BooleanAnalysis.GaussianSpace.measurableEmbedding_gaussianTorusMap}
\label{BooleanAnalysis.GaussianSpace.measurableEmbedding_gaussianTorusMap}
\leanok
\uses{BooleanAnalysis.GaussianSpace.continuous_gaussianTorusMap, BooleanAnalysis.GaussianSpace.gaussianTorusMap, BooleanAnalysis.GaussianSpace.gaussianTorusMap_injective}
\difficulty{1}
For each $n\ge0$, map $x\in\mathbb R^n$ to the torus with coordinates indexed by
$(j,b)\in\{1,\ldots,n\}\times\{0,1\}$, whose $(j,0)$ coordinate is $x_j$ modulo $\mathbb
Z$ and whose $(j,1)$ coordinate is $\sqrt2 x_j$ modulo $\mathbb Z$. This map is a
measurable embedding.
\end{lemma}

\begin{lemma}[Fourier character of the Gaussian torus map]
\lean{BooleanAnalysis.GaussianSpace.mFourier_gaussianTorusMap}
\label{BooleanAnalysis.GaussianSpace.mFourier_gaussianTorusMap}
\leanok
\uses{BooleanAnalysis.GaussianSpace.gaussianTorusMap}
\difficulty{3}
Let $n\ge 0$, let $m:\{1,\ldots,n\}\times\{0,1\}\to\mathbb Z$, and let $x\in\mathbb
R^n$. The Fourier character of the Gaussian torus image of $x$ is
$\exp\!\left(i\sum_{j=1}^n2\pi\bigl(m_{j,0}+\sqrt2\,m_{j,1}\bigr)x_j\right)$.
\end{lemma}

\begin{lemma}[Density of torus Fourier characters in $L^p$]
\lean{BooleanAnalysis.GaussianSpace.span_mFourier_toLp_dense}
\label{BooleanAnalysis.GaussianSpace.span_mFourier_toLp_dense}
\leanok
\difficulty{1}
Let $\iota$ be a finite index set, let $\mu$ be a finite measure on the torus with
coordinates indexed by $\iota$, and let $p$ be a finite real exponent with $1\le p$. The
complex linear span of all torus Fourier characters is dense in $L^p(\mu;\mathbb C)$.
\end{lemma}

\begin{lemma}[Surjectivity of Gaussian torus pullback]
\lean{BooleanAnalysis.GaussianSpace.surjective_gaussianTorus_pullback}
\label{BooleanAnalysis.GaussianSpace.surjective_gaussianTorus_pullback}
\leanok
\uses{BooleanAnalysis.GaussianSpace.continuous_gaussianTorusMap, BooleanAnalysis.GaussianSpace.gaussianTorusMap, BooleanAnalysis.GaussianSpace.measurableEmbedding_gaussianTorusMap}
\difficulty{5}
Let $n\ge0$, let $\mu$ be any measure on $\mathbb R^n$, let $\mathbb K$ be either
$\mathbb R$ or $\mathbb C$, and let $p\in[1,\infty]$. Define $T:\mathbb R^n\to(\mathbb
R/\mathbb Z)^{\{1,\ldots,n\}\times\{0,1\}}$ by $T(x)_{i,0}=x_i\pmod{\mathbb Z}$ and
$T(x)_{i,1}=\sqrt2\,x_i\pmod{\mathbb Z}$. The pullback map $f\mapsto f\circ T$ from
$L^p(T_\#\mu)$ to $L^p(\mu)$ is surjective.
\end{lemma}

\begin{definition}[Fourier character in Gaussian-space Lp]
\lean{BooleanAnalysis.GaussianSpace.gaussianTorusFourierLp}
\label{BooleanAnalysis.GaussianSpace.gaussianTorusFourierLp}
\leanok
\uses{BooleanAnalysis.GaussianSpace.continuous_gaussianTorusMap, BooleanAnalysis.GaussianSpace.gaussianTorusMap}
Let $n$ be a natural number, let $\bbc$ be a finite measure on $\bbr^n$, let
$p\in[1,\infty]$, and let $m$ assign an integer to each pair in
$\{1,\ldots,n\}\times\{0,1\}$. The Gaussian-torus Fourier function associated with $m$,
composed with the Gaussian torus map, defines an element of $L^p(\mu;\bbc)$.
\end{definition}

\begin{lemma}[Gaussian torus Fourier representative]
\lean{BooleanAnalysis.GaussianSpace.coeFn_gaussianTorusFourierLp}
\label{BooleanAnalysis.GaussianSpace.coeFn_gaussianTorusFourierLp}
\leanok
\uses{BooleanAnalysis.GaussianSpace.continuous_gaussianTorusMap, BooleanAnalysis.GaussianSpace.gaussianTorusFourierLp, BooleanAnalysis.GaussianSpace.gaussianTorusMap}
\difficulty{3}
Let $n$ be a natural number, let $\mu$ be a finite measure on $\mathbb R^n$, let $p\ge
1$ be an extended nonnegative real number, and let
$m:(\mathrm{Fin}\,n\times\{0,1\})\to\mathbb Z$. The $L^p(\mu)$ function obtained by
pulling back the torus Fourier character indexed by $m$ along the Gaussian torus map
agrees $\mu$-almost everywhere with the function that evaluates that character at the
image of each point under the map.
\end{lemma}

\begin{lemma}[Density of Gaussian torus Fourier functions]
\lean{BooleanAnalysis.GaussianSpace.span_gaussianTorusFourier_dense}
\label{BooleanAnalysis.GaussianSpace.span_gaussianTorusFourier_dense}
\leanok
\uses{BooleanAnalysis.GaussianSpace.continuous_gaussianTorusMap, BooleanAnalysis.GaussianSpace.gaussianTorusFourierLp, BooleanAnalysis.GaussianSpace.gaussianTorusMap, BooleanAnalysis.GaussianSpace.span_mFourier_toLp_dense, BooleanAnalysis.GaussianSpace.surjective_gaussianTorus_pullback}
\difficulty{4}
Let $n$ be a nonnegative integer, let $\mu$ be a finite measure on $\mathbb R^n$, and
let $p$ be a finite real exponent with $1\le p$. Map $x\in\mathbb R^n$ to the torus with
coordinates indexed by pairs
$(j,\varepsilon)\in\{0,\ldots,n-1\}\times\{\mathrm{false},\mathrm{true}\}$, using the
circle class of $x_j$ when $\varepsilon=\mathrm{false}$ and of $\sqrt2\,x_j$ when
$\varepsilon=\mathrm{true}$. The complex linear span of the pullbacks, under this map,
of all torus Fourier characters is dense in $L^p(\mu;\mathbb C)$.
\end{lemma}

\begin{definition}[Gaussian torus frequency]
\lean{BooleanAnalysis.GaussianSpace.gaussianTorusFrequency}
\label{BooleanAnalysis.GaussianSpace.gaussianTorusFrequency}
\leanok
For a natural number $n$, an integer-valued assignment $m$ on
$\mathrm{Fin}(n)\times\{0,1\}$, and $i\in\mathrm{Fin}(n)$, define the real frequency
$\omega_i=2\pi\bigl(m(i,0)+\sqrt{2}\,m(i,1)\bigr)$.
\end{definition}
